\documentclass[11pt,a4paper]{article}
\usepackage[margin=2cm]{geometry}
\usepackage{amsmath,amssymb,bm,booktabs,microtype}
\usepackage[colorlinks=true,linkcolor=blue,citecolor=blue]{hyperref}
\setlength{\parskip}{2pt}
\title{\vspace{-1.6cm}\textbf{From \texttt{JPCB\_tandem.tex} to
\texttt{JPCB\_tandem\_round\_2.tex}:\\ what changed, and how the CD simulation works}}
\author{}
\date{\vspace{-1.2cm}}

\begin{document}
\maketitle
\vspace{-1.0cm}

\section*{1. What changed between the two manuscripts}

The two documents reach opposite conclusions from the same trajectory. The cause
is a single unit-conversion error, and the consequence is a different physical
story.

\begin{table}[h]
\centering\small
\begin{tabular}{lll}
\toprule
& \textbf{Round 1} (\texttt{JPCB\_tandem.tex}) & \textbf{Round 2} (\texttt{JPCB\_tandem\_round\_2.tex}) \\
\midrule
Title & unchanged & unchanged (minimal revision of the same paper) \\
Coupling $J$ & $117.2 \pm 5.5~\mathrm{cm^{-1}}$ & $\bm{32.82 \pm 1.55~\mathrm{cm^{-1}}}$ \\
Point dipole & $27.6 \pm 1.3~\mathrm{cm^{-1}}$ & $27.6 \pm 1.3~\mathrm{cm^{-1}}$ (unchanged) \\
TDC/PDA & $4.24$ ``near-field multipolar'' & $\bm{1.187 \pm 0.011}$ \\
Detuning $\Delta$ & not computed; degeneracy assumed & $\langle|\Delta|\rangle = 576~\mathrm{cm^{-1}}$ ($n=95$) \\
Eigenstates & delocalised exciton & localised, mean mixing $0.18$ \\
Agreement with CD & $J$ within $10\%$ of $130.79$ & couplet separation is not a gap estimator \\
Experimental check & none available & 2PA polarisation ratio, $32$--$40~\mathrm{cm^{-1}}$ unscreened \\
Dynamics & single illustrative trajectory & two extreme ensemble frames, $\Delta = 12.9$ and $-1825$ \\
\bottomrule
\end{tabular}
\end{table}

\paragraph{The error.} The Coulomb double sum is evaluated with grid
coordinates in \AA{}, so the reciprocal distance must be converted with
$1~\text{\AA}^{-1} = 0.529177\,a_0^{-1}$. Round~1 applied the \emph{coordinate}
factor $1.8897$ instead, inflating every transition-density coupling by
$1.8897^{2} = 3.5711$ while leaving the point-dipole values --- evaluated in
bohr throughout --- untouched. Exactly:
$117.2/3.5711 = 32.82~\mathrm{cm^{-1}}$. The earlier arXiv preprint carried a
\emph{different} version of the same mistake (no conversion at all, inflation
$1.8897$), so its $74.38~\mathrm{cm^{-1}}$ must not be rescaled by $3.5711$.

\paragraph{Why it changes the physics.} At $4.24\times$ the point-dipole value,
the coupling looked large enough to rival the CD splitting, and delocalisation
followed. At $1.19\times$ it is an ordinary dipole--quadrupole correction, the
coupling is $32.8~\mathrm{cm^{-1}}$, and it is an order of magnitude
\emph{below} the site-energy detuning. Round~2 therefore computes $\Delta$,
which round~1 never did, and finds the eigenstates localised. Independent
support has since appeared: a referee's Q-Chem calculation
($6.02$ and $5.36$~meV, ratio $1.12$) and our matched ORCA control
($6.087$ and $5.361$~meV, ratio $1.14$); and Cusick \emph{et al.} (\emph{JPCA}
\textbf{2026}, 130, 5471), from the laboratories that reported the large
couplings, now compute $26.9$--$36.8~\mathrm{cm^{-1}}$ for this dimer.

\section*{2. How the CD simulation is obtained}

\paragraph{Per-frame Hamiltonian.} Each MD snapshot supplies a two-site
Hamiltonian $H = \begin{pmatrix} E_A & J \\ J & E_B\end{pmatrix}$, with $E_A$,
$E_B$ from QM/MM TDDFT (CAM-B3LYP/6-311+G**, full-system embedding) and $J$ from
the integrated DLPNO-STEOM-CCSD transition densities. Diagonalising gives
\begin{equation}
\Omega = \sqrt{\Delta^{2}+4J^{2}}, \qquad
\tan 2\theta = \frac{2J}{\Delta}, \qquad \Delta = E_A - E_B ,
\end{equation}
with eigenstates $|\pm\rangle$ at $\bar{E} \pm \Omega/2$.

\paragraph{Why a two-state block is enough.} The basis also contains
$|0\rangle$ and $|e_Ae_B\rangle$, and the first is not weakly coupled:
$\langle 0|\hat V|1\rangle$ pairs one chromophore's \emph{transition} density
with the other's \emph{permanent} charge, and since the chromophore is anionic
that is a $-1\,e$ monopole, giving an $R^{-2}$ term of $292~\mathrm{cm^{-1}}$
unscreened --- nine times $J$. It is negligible for a reason central to this
paper: unlike $J$, screened by $\varepsilon_{\mathrm{opt}}$ because it couples
transition densities at optical frequency, this term acts against a static
charge the environment has equilibrated to before absorption. With the optical
gap as denominator the second-order correction to $J$ is $1.42~\mathrm{cm^{-1}}$
at $\varepsilon_{\mathrm{opt}}$, $0.28$ at $\varepsilon=4$ and $0.001$ at
$\varepsilon_s=78$; the diagonal effect of the same monopole is already inside
the QM/MM site energies. $|e_Ae_B\rangle$ lies $580\,J$ above and is irrelevant
here, but it is what makes two-photon absorption a distinct observable, so
two-photon and antibunching data need the full $4\times4$ block.

\paragraph{Rotational strengths.} The exciton-chirality rotational strength of
the two bands is equal and opposite,
\begin{equation}
R_{\pm} = \mp\,\frac{\pi}{2}\,\bar\nu\,\sin(2\theta)\; T,
\qquad
T = \bm{R}_{AB}\cdot(\bm{\mu}_A \times \bm{\mu}_B),
\label{eq:R}
\end{equation}
where $T$ is the chirality triple product, computed per frame from the same
placed transition densities that give $J$, so that the sign of $J$ and the sign
of $T$ are guaranteed mutually consistent. \textbf{The key structural feature of
Eq.~\eqref{eq:R} is the factor $\sin 2\theta = 2J/\Omega$}: the couplet
\emph{amplitude} carries the coupling, and it is suppressed by $2J/\Omega =
0.119$ at the computed detuning. The couplet \emph{separation} is $\Omega$,
which is dominated by $\Delta$. Amplitude and width therefore measure different
things, and only the amplitude measures $J$.

\paragraph{Lineshape and ensemble.} Each band is broadened with a Voigt profile
whose homogeneous half-width is $1/(2\pi c\,T_2^{*})$ and whose inhomogeneous
width comes from the frame-to-frame spread; the spectrum is then averaged over
all $1000$ frames. Because $\Omega$ enters nonlinearly, this ensemble average is
\emph{not} the spectrum of the mean Hamiltonian: the rotational-strength-weighted
mean of $\Omega$ lies close to its harmonic mean ($295.9~\mathrm{cm^{-1}}$)
rather than its arithmetic mean ($611.2~\mathrm{cm^{-1}}$). Vibronic structure is
added in the one-particle Holstein approximation ($\omega \approx
1450~\mathrm{cm^{-1}}$, $S = 0.35$--$0.5$), which supplies the third band.

\paragraph{Comparison protocol.} Simulated spectra are put through Nguyen
\emph{et al.}'s own analysis rather than compared to it directly: the same
$\Delta[\theta] = [\theta]_{\mathrm{TDX}} - [\theta]_{\mathrm{TD}}$ subtraction
order (which carries \emph{minus} the dimer's own couplet, and is the origin of
an apparent handedness disagreement), and the same constrained three-Lorentzian
fit with the couplet centres pinned symmetrically about the TDX absorption peak.
The absolute scale uses the published $\varepsilon_{516}$ and $25~\mu$M
concentration.

\section*{3. Which experiments the simulation agrees with}

\begin{table}[h]
\centering\small
\begin{tabular}{llll}
\toprule
Observable & Experiment & This work & Status \\
\midrule
CD couplet separation & $14.6 \pm 0.3$~nm (Kim 2019) & lineshape-dominated; not a gap estimator & \textbf{not comparable} \\
Apparent coupling & $130.79~\mathrm{cm^{-1}}$ (constrained fit) & $105$--$159~\mathrm{cm^{-1}}$ & agrees \\
Couplet handedness & negative (after subtraction order) & negative in every frame & agrees \\
Third CD band & $481.4$~nm & vibronic, $470$--$495$~nm & agrees \\
Absorption red shift & $-35~\mathrm{cm^{-1}}$ & $-12$ to $-16$ electrostatic $+$ excitonic & agrees \\
Radiative shortening & $57 \pm 4$~ps & $42$--$58$~ps at $J \approx 33$ & agrees \\
Limiting anisotropy & $0.52 \to 0.30$ & requires $\alpha \approx 131^\circ$ vs computed $110^\circ$ & \textbf{open} \\
\bottomrule
\end{tabular}
\end{table}

The first row is the one that changed most between versions, and it is now a
negative result rather than a positive one. Kim's splitting has been read
throughout the literature as $2|J|$, giving $J = 274~\mathrm{cm^{-1}}$; at the
crystallographic $25.4$~\AA{} separation that requires a transition dipole of
$\approx 29$~D or a separation of $\approx 10$~\AA. An earlier version of this
work proposed reading it instead as the detuned exciton gap $\Omega$, which
appeared to reproduce $14.6$~nm with nothing adjusted. Direct simulation of the
couplet withdraws that claim: propagating $\Omega = 553.9~\mathrm{cm^{-1}}$
through the lineshape of the published constrained fit gives extrema separated
by $652~\mathrm{cm^{-1}}$, not the observed $548$. On that lineshape the
separation is a non-monotonic, two-valued and linewidth-dominated function of
the gap --- a gap of \emph{zero} still gives $690~\mathrm{cm^{-1}}$ --- so it
estimates the gap in neither direction, and the observable was never comparable
with a computed $J$ or $\Omega$ in the first place. The detuning survives as the
explanation of the couplet \emph{amplitude} and of the localisation of the
eigenstates, and the same artefact is visible and correctly handled in the
two-photon $\Omega$ spectrum of Cusick et al.

The one row that does not yet agree is the limiting anisotropy, which is
independently corroborated ($49^\circ$ for the eYFP dimer by one-photon
anisotropy, Kinoshita 2026) and which we therefore treat as an open structural
question rather than a resolved one.

\end{document}
